\documentclass[a4paper,11pt,twoside]{article}
\usepackage[utf8]{inputenc}	% Text coding
\usepackage[T1]{fontenc}
\usepackage{lmodern}
\usepackage[czech]{babel}
\usepackage{epsfig}
\usepackage{amsfonts,amsmath,amssymb}
\usepackage{graphicx}
\usepackage[unicode]{hyperref}
%\usepackage{indentfirst}
\usepackage{fancyhdr}
\usepackage{xifthen}
\usepackage{amsthm,thmtools}
\usepackage{bold-extra}
\usepackage[dvipsnames]{xcolor}
\usepackage[subrefformat=simple,labelformat=simple]{subcaption} % Instead of subfigure
\usepackage{listings}
\usepackage{comment}
\usepackage{titlesec}
\usepackage{underscore}
\usepackage{makecell}       % Šířky čar v tabulkách

% Page size
\addtolength{\topmargin}{-1.5cm} %\addtolength{\textheight}{-10cm}
\addtolength{\textwidth}{4cm} \addtolength{\textheight}{4cm} % Width and height of the text
\addtolength{\voffset}{-0.5cm} % Top margin
\addtolength{\hoffset}{-2cm}
\setlength{\headheight}{15pt}

\DeclareMathOperator{\e}{e}

\def\vector#1{\boldsymbol{#1}}								% Vector
\renewcommand{\d}{\mathrm{d}}
\newcommand{\derivative}[3][]{\ifthenelse{\isempty{#1}}	    % Normal derivative
	{\frac{\d{#2}}{\d{#3}}}
	{\frac{\d^{#1}{#2}}{\d{#3}^{#1}}}
}
\newcommand{\im}{\mathrm{i}}

\def\makematrix#1{\begin{pmatrix}#1\end{pmatrix}}       % Matrix
\def\abs#1{\left|#1\right|}
\def\probability#1{\mathrm{Pr}\left[#1\right]}
\def\expectation#1{\mathrm{E}\left[#1\right]}
\def\dispersion#1{\sigma_{#1}^{2}}
\def\c{,\!}

\def\code#1{\textnormal{\texttt{#1}}}
\def\file#1{\textnormal{\textbf{\texttt{#1}}}}
\def\ghfile#1#2{\textnormal{\textbf{\texttt{\href{https://github.com/PavelStransky/PCInPhysics2021/blob/main/#1#2}{#2}}}}}

\def\abbreviation#1{\textnormal{\textsc{#1}}}

\def\ket#1{\left|{#1}\right\rangle}
\def\bra#1{\left\langle{#1}\right|}
\def\braket#1#2{\left\langle{#1}\middle|{#2}\right\rangle}

\begin{document}

\section*{Domácí úkol na 12.5.2025}
\subsection*{Wolfram Mathematica}

V programu Wolfram Mathematica vyřešte a vhodně naformátujte následující úlohy.

\begin{enumerate}
    \item
    Spočítejte limitu
    \begin{equation*}
        \lim_{x\rightarrow 0}\frac{\sqrt[3]{\sin{x}+x^3}-x}{x^5}.
    \end{equation*}

    \item
    Spočítejte určitý integrál
    \begin{equation*}
        I=\int_{0}^{\infty}\frac{\ln{x}}{\left(x^2+1\right)^2}\,\d x.
    \end{equation*}
\end{enumerate}

\subsubsection*{Dvouhladinový kvantový systém}

Uvažujte dvouhladinový systém (například spin) popsaný Hamiltoniánem (maticí $2\times2$)
\begin{equation*}
    H=\begin{pmatrix}
        \Delta & \Omega \\
        \Omega & \Delta
    \end{pmatrix},
\end{equation*}
kde $\Delta$ je energetický rozdíl mezi dvěma hladinami a $\Omega$ udává intenzitu interakce.

\begin{enumerate}
    \setcounter{enumi}{2}
    \item 
    Nalezněte vlastní hodnoty (energie) a normalizované vlastní vektory (odpovídající stavy) Hamiltoniánu $H$.
    
    \item 
    Vykreslete vlastní hodnoty jako funkci $\Delta$ pro dané $\Omega=1$.

    \item
    Vykreslete komponenty vlastních vektorů jako funkci $\Delta$ pro dané $\Omega=1$.

    \item
    Časový vývoj stavu systému je dán řešením Schrödingerovy rovnice, které v tomto případě lze explicitně nalézt jako
    \begin{equation*}
        \ket{\psi(t)}=\exp\left(-\im H t\right)\ket{\psi_{0}},
    \end{equation*}    
    kde
    \begin{equation*}
        \ket{\psi_{0}}=\begin{pmatrix}1 \\ 0\end{pmatrix}
    \end{equation*}
    popisuje stav systému v čase $t=0$.
    Napište funkci, která spočítá pravděpodobnost naměření stavu $\ket{\psi_{0}}$ v čase $t$ pro libovolné hodnoty $\Delta$ a $\Omega$.
    Pravděpodobnost se vypočítá jako
    \begin{equation*}
        p(t;\Delta,\Omega)=\abs{\braket{\psi_{0}}{\psi(t)}}^2.
    \end{equation*}
    Pravděpodobnost vykreslete jako funkci $t$ pro nějaké zvolené $\Omega$ a $\Delta$.

    Symbol $\braket{\phi}{\psi}$ značí skalární součin a $\bra{\psi_{0}}$ je duální vektor k $\ket{\psi_{0}}$,
    \begin{equation*}
        \bra{\psi_{0}}=\begin{pmatrix}1 & 0\end{pmatrix}.
    \end{equation*}

\end{enumerate}

\end{document}